\chapter{The Toolkit}
\label{ch:toolkit}

``So what are the sixteen modes?''\footnote{The Pali: \textit{``Tattha katame soḷasa hārā? Desanā vicayo yutti padaṭṭhāno lakkhaṇo catubyūho āvaṭṭo vibhatti parivattano vevacano paññatti otaraṇo sodhano adhiṭṭhāno parikkhāro samāropano iti.''} Mahākaccāyana asks and answers in one breath. The word \textit{hāra} (which we translate ``mode'') literally means ``carrying, conveying'' from the root \textit{har}. Ñāṇamoli rendered it ``conveying'' --- technically right but gives English readers nothing to hold onto. We chose ``mode'' because these are genuinely \textit{modes of reading}: different lenses you apply to the same text to see different things.}

Just like that, no preamble. Mahākaccāyana asks the question and answers it in a single breath. Here they are --- your sixteen tools for reading any sutta the Buddha ever taught:

\begin{enumerate}[leftmargin=*, itemsep=8pt]
\item \textbf{Teaching} --- How does the Buddha structure a lesson? There's a pattern: here's the appeal, here's the danger, here's the way out, here's the result, here's how to do it, and --- do it. Every time.\footnote{\textit{Desanāhāra} --- the ``teaching mode.'' From \textit{deseti}, ``to point out.'' The five-part structure is: gratification (\textit{assāda}), danger (\textit{ādīnava}), escape (\textit{nissaraṇa}), fruit (\textit{phala}), and method (\textit{upāya}), plus his direct instruction (\textit{āṇatti}) to the practitioner. This pattern appears across hundreds of suttas. Once you see it, you can't unsee it.}

\item \textbf{Investigation} --- What's the question being asked? What's the answer? What's the summary? Break any sutta into these three and you can see its spine.\footnote{\textit{Vicayahāra} --- from \textit{vicināti}, ``to investigate, to sort through.'' The three components: what question is being asked (\textit{pucchita}), what answer is given (\textit{vissajjita}), and what summary verse accompanies the teaching (\textit{anugīti}).}

\item \textbf{Fitness} --- Does this interpretation actually fit? This is the quality-control mode. It checks whether your reading of a passage is consistent with everything else the Buddha taught. If it doesn't fit, you've got it wrong.\footnote{\textit{Yuttihāra} --- from \textit{yujjati}, ``to be fitting, to connect.'' This is a meta-mode: it examines whether a particular reading or interpretation is appropriate and consistent with the overall teaching. It's the Netti's quality control system. You could think of it as peer review for sutta interpretation.}

\item \textbf{Basis} --- What's the foundation? When the Buddha talks about a quality --- say, mindfulness --- what does it rest on? What supports it? Trace any teaching down to what it stands on.\footnote{\textit{Padaṭṭhānahāra} --- literally ``the footing,'' from \textit{pada} (foot, basis) + \textit{ṭhāna} (standing place). The Pali: \textit{``Dhammaṃ deseti jino, tassa ca dhammassa yaṃ padaṭṭhānaṃ''} --- ``The Victor teaches a quality; this mode finds what that quality stands on.''}

\item \textbf{Characteristic} --- What's the shared signature? When the Buddha describes one thing, which other things share its defining mark? This is how a teaching about one topic secretly covers fifty topics at once.\footnote{\textit{Lakkhaṇahāra} --- ``mark, characteristic, sign.'' The verse summary: \textit{``Vuttamhi ekadhamme, ye dhammā ekalakkhaṇā keci; / vuttā bhavanti sabbe''} --- ``When one thing with a certain characteristic is stated, all things sharing that characteristic are [implicitly] stated.'' This is a remarkably modern insight about the logic of universals.}

\item \textbf{The Fourfold Analysis} --- Four lenses on any passage: What does the grammar actually mean? What's the intention behind the words? What's the context? And how does this passage connect to what comes before and after it?\footnote{\textit{Catubyūhahāra} --- literally ``four-array.'' The four: grammatical intention (\textit{nirutti}), underlying purpose (\textit{adhippāya}), phrasing and background (\textit{byañjana} and \textit{nidāna}), and the connection between what comes before and after (\textit{pubbāparānusandhi}). This is essentially ancient discourse analysis.}

\item \textbf{Turning Back} --- Follow the chain \textit{and} flip it. Find what a teaching rests on, then find what \textit{that} rests on. Then turn to the opposite. If you're looking at greed, turn to non-greed. If you're looking at the problem, turn to the solution.\footnote{\textit{Āvaṭṭahāra} --- from \textit{āvaṭṭati}, ``to turn back, to revolve.'' The verse: \textit{``Ekamhi padaṭṭhāne, pariyesati sesakaṃ padaṭṭhānaṃ; / āvaṭṭati paṭipakkhe''} --- ``From one basis, it searches for the remaining basis; it turns to the opposite.''}

\item \textbf{Sorting} --- Which parts of this teaching apply to everyone? Which parts are specific to a particular level of practice? Not everything the Buddha said applies to every person at every stage.\footnote{\textit{Vibhattihāra} --- from \textit{vibhajati}, ``to divide, to sort.'' Sorts teachings into shared (\textit{sādhāraṇa}) and specific (\textit{asādhāraṇa}) applications.}

\item \textbf{Reversal} --- Flip the script. When the Buddha says ``develop this,'' what's being abandoned? When he says ``abandon this,'' what's being developed? Every positive instruction has a negative shadow, and vice versa.\footnote{\textit{Parivattanahāra} --- from \textit{parivattati}, ``to turn around, to reverse.'' The verse: \textit{``Kusalākusale dhamme, niddiṭṭhe bhāvite pahīne ca; / parivattati paṭipakkhe''} --- ``Among wholesome and unwholesome qualities, described, developed, abandoned --- it turns to the opposite.''}

\item \textbf{Synonyms} --- Same idea, different words. The Buddha uses multiple terms for the same thing across different suttas. Spot the synonyms and the whole Canon starts to cross-reference itself.\footnote{\textit{Vevacanahāra} --- ``synonym, alternative expression.'' The Buddha deliberately uses different words for the same concept in different suttas. For example, \textit{taṇhā} (craving), \textit{nandī} (delight), and \textit{upādāna} (clinging) can all point to the same fundamental process, viewed from different angles. The verse: \textit{``Vevacanāni bahūni tu, sutte vuttāni ekadhammassa''} --- ``Many synonyms for one and the same thing are stated in the suttas.''}

\item \textbf{Designation} --- Different names for the same reality. The Buddha teaches the same truth through many different frameworks --- the five aggregates, the twelve sense bases, the eighteen elements. These aren't competing systems. They're different windows into the same room.\footnote{\textit{Paññattihāra} --- from \textit{paññāpeti}, ``to make known, to designate.'' The verse: \textit{``Ekaṃ bhagavā dhammaṃ, paññattīhi vividhāhi deseti''} --- ``The Blessed One teaches one reality through various designations.'' The diversity of frameworks is pedagogical, not ontological.}

\item \textbf{Entry Point} --- Any teaching can be mapped onto the standard frameworks: dependent origination, the aggregates, the sense bases, the elements. These are the doors into deeper understanding.\footnote{\textit{Otaraṇahāra} --- from \textit{otarati}, ``to descend into, to enter.'' The verse: \textit{``Yo ca paṭiccuppādo, indriyakhandhā ca dhātu āyatanā; / etehi otarati yo, otaraṇo nāma so hāro''} --- ``Dependent origination, faculties, aggregates, elements, sense bases --- one who enters through these, that mode is called Entry Point.''}

\item \textbf{Clearing} --- Does the answer actually match the question? When a sutta poses a question in verse and answers it in prose, this mode checks whether the answer really covers what was asked. Think of it as debugging.\footnote{\textit{Sodhanahāra} --- from \textit{sodheti}, ``to purify, to clear, to verify.'' It checks whether an answer is ``pure'' (\textit{suddha} --- fully adequate) or ``impure'' (\textit{asuddha} --- only partial). The Pali: \textit{``Suddhāsuddhaparikkhā, hāro so sodhano nāma.''}}

\item \textbf{Determination} --- Is this the same thing or a different thing? When the Buddha describes two qualities, are they the same reality described differently, or genuinely distinct? This mode sorts identity from difference.\footnote{\textit{Adhiṭṭhānahāra} --- from \textit{adhitiṭṭhati}, ``to stand upon, to determine.'' Teachings describe things that are ``the same in nature'' (\textit{ekattatā}) or ``different in nature'' (\textit{vemattatā}). The verse: \textit{``Ekattatāya dhammā, yepi ca vemattatāya niddiṭṭhā; / tena vikappayitabbā.''}}

\item \textbf{Requisites} --- What do you need in order to get what you need? This mode traces the causal chain: which quality is the prerequisite for which other quality? It maps the supply chain of the path.\footnote{\textit{Parikkhārahāra} --- ``requisite, equipment.'' The verse: \textit{``Ye dhammā yaṃ dhammaṃ, janayanti paccayā paramparato; / hetumavakaḍḍhayitvā''} --- ``Which qualities generate which quality, as conditions in sequence --- having drawn out the cause.'' This mode traces the causal DNA of the path.}

\item \textbf{Integration} --- Put it all together. After you've broken things apart, this mode brings them back together. Teachings that share the same root, the same purpose, the same direction --- collect them. Synthesize. See the whole.\footnote{\textit{Samāropanahāra} --- from \textit{samāropeti}, ``to make ascend together, to integrate.'' The verse: \textit{``Ye dhammā yaṃ mūlā, ye cekatthā pakāsitā muninā; / te samaropayitabbā''} --- ``Those teachings that share the same root, that are proclaimed by the Sage as having one purpose --- those should be integrated.'' This is the capstone mode: analysis gives way to synthesis.}
\end{enumerate}

\begin{versebox}
Teaching, Investigation, Fitness, Basis, Characteristic,\\
Fourfold Analysis, Turning Back, Sorting, Reversal,\\
Synonyms, Designation, Entry Point, Clearing,\\
Determination, Requisites, Integration --- sixteen.
\stanza
These sixteen modes are stated, each distinct in purpose.\\
And from these very modes, expanded, the methods are distributed.\footnote{Summary verse (\textit{Tassānugīti}): \textit{``Desanā vicayo yutti, padaṭṭhāno ca lakkhaṇo; / catubyūho ca āvaṭṭo, vibhatti parivattano. / Vevacano ca paññatti, otaraṇo ca sodhano; / adhiṭṭhāno parikkhāro, samāropano soḷaso. // Ete soḷasa hārā, pakittitā atthato asaṃkiṇṇā; / etesañceva bhavati, vitthāratayā nayavibhattīti.''}}
\end{versebox}

Notice the last line: ``from these very modes, expanded, the methods are distributed.'' The modes and the methods aren't separate systems. The five methods \textit{emerge from} the sixteen modes. They're what happens when you zoom out from individual reading techniques to the big-picture patterns. Let's see them.

\ornament

\section*{The Five Methods}

``And what are the five methods?''\footnote{The Pali: \textit{``Tattha katame pañca nayā? Nandiyāvaṭṭo tipukkhalo sīhavikkīḷito disālocano aṅkuso iti.''} The word \textit{naya} means ``method,'' ``way,'' or ``leading'' --- from the same root as the title \textit{Netti} itself.} Again, straight to it:

\begin{enumerate}[leftmargin=*, itemsep=8pt]
\item \textbf{The Joyful Round} --- Map any teaching onto craving and ignorance on one side, calm and insight on the other, and connect them through the Four Noble Truths.\footnote{\textit{Nandiyāvaṭṭa} --- literally ``the turning of delight.'' The name evokes a cycle: craving and ignorance (the unwholesome pair) are countered by calm and insight (the wholesome pair), mapped onto the Four Noble Truths. The verse: \textit{``Taṇhañca avijjampi ca, samathena vipassanā yo neti; / saccehi yojayitvā, ayaṃ nayo nandiyāvaṭṭo.''}}

\item \textbf{The Triple Tip} --- Take any unwholesome root and its wholesome opposite, and see how the teaching connects them. Three roots, three antidotes, three points on every teaching.\footnote{\textit{Tipukkhala} --- literally ``three-pointed,'' like a trident. The three ``tips'' are: the unwholesome quality, its wholesome counterpart, and the truth that connects them. The verse: \textit{``Yo akusale samūlehi, neti kusale ca kusalamūlehi; / bhūtaṃ tathaṃ avitathaṃ, tipukkhalaṃ taṃ nayaṃ āhu.''}}

\item \textbf{The Lion's Play} --- Like a lion surveying every direction. Take the four distortions of perception --- seeing permanence where there's impermanence, pleasure where there's suffering, self where there's non-self, beauty where there's the unattractive --- and their corrections. Turn in every direction through the teaching.\footnote{\textit{Sīhavikkīḷita} --- ``the lion's sport.'' A lion turns in every direction to survey its territory before acting. This method takes the four perceptual distortions (\textit{vipallāsa}) and their corrections through the four foundations of mindfulness. The image is of sovereign confidence: the reader-as-lion surveys the entire terrain of a sutta.}

\item \textbf{Scanning the Horizon} --- Survey the whole landscape. Before you dive in, scan the entire teaching and note every wholesome and unwholesome quality that appears. Get the lay of the land.\footnote{\textit{Disālocana} --- ``looking in all directions,'' from \textit{disā} (direction) + \textit{ālocana} (examining). The reconnaissance method. The verse calls it \textit{``nayamuttamaṃ''} --- ``the supreme method.'' Before doing the detailed work, look everywhere and note what's there.}

\item \textbf{The Hook} --- After scanning, hook out what you've found and sort it. Gather all the wholesome qualities together, all the unwholesome qualities together. Pull the essential structure out of the mass of words.\footnote{\textit{Aṅkusa} --- a goad or hook, the kind used by an elephant trainer. Having ``scanned the horizon,'' this method ``hooks'' out the essential elements --- the way a mahout guides an elephant. The summary verse: \textit{``Oloketvā disalocanena, ukkhipiya yaṃ samāneti; / sabbe kusalākusale, ayaṃ nayo aṅkuso nāma.''} --- ``Having surveyed with the Horizon Scan, lifting up and gathering together all the wholesome and unwholesome --- this method is called the Hook.''}
\end{enumerate}

\begin{versebox}
First the Joyful Round, and then the Triple Tip,\\
the Lion's Play comes third --- the method's hallmark.
\stanza
Scanning the Horizon fourth, supreme of all;\\
fifth, the Hook --- and that's all five methods told.
\end{versebox}

Notice how the methods build on each other. You start with the Joyful Round (the simplest mapping: craving and ignorance versus calm and insight, through the Four Truths), get more specific with the Triple Tip (the three roots and their opposites), go wider with the Lion's Play (the four distortions and their corrections), then survey everything (Scanning the Horizon), and finally gather and synthesize (the Hook).

It's a zoom sequence. Close-up, medium, wide angle, panorama, and then you collect your shots into the final picture.

\ornament

\section*{The Eighteen Root Terms}

And here's the bedrock.\footnote{The Pali: \textit{``Tattha katamāni aṭṭhārasa mūlapadāni? Nava padāni kusalāni nava padāni akusalāni.''} --- ``And what are the eighteen root terms? Nine wholesome terms and nine unwholesome terms.'' The word \textit{mūlapada} means literally ``root-word'' or ``root-term'' --- these are the irreducible elements to which everything else connects.} Everything the Buddha ever taught, according to this manual, connects back to eighteen root terms. Nine unwholesome:

\begin{enumerate}[itemsep=2pt]
\item \textbf{Craving} --- the raw pull toward experience
\item \textbf{Ignorance} --- not seeing things as they are
\item \textbf{Greed} --- wanting what's in front of you
\item \textbf{Hatred} --- pushing away what's in front of you
\item \textbf{Delusion} --- being confused about what's in front of you
\item \textbf{The perception ``this is beautiful''} --- seeing attractiveness where it doesn't ultimately hold
\item \textbf{The perception ``this is pleasant''} --- seeing pleasure where there's really suffering
\item \textbf{The perception ``this is permanent''} --- seeing stability where everything is changing
\item \textbf{The perception ``this is me''} --- seeing a self where there's only process
\end{enumerate}

And nine wholesome:

\begin{enumerate}[itemsep=2pt]
\item \textbf{Calm} --- settling the mind
\item \textbf{Insight} --- seeing clearly
\item \textbf{Non-greed} --- letting go
\item \textbf{Non-hatred} --- goodwill
\item \textbf{Non-delusion} --- understanding
\item \textbf{The perception ``this is not beautiful''} --- seeing through surface attractiveness
\item \textbf{The perception ``this is suffering''} --- recognizing the unsatisfactoriness of conditioned experience
\item \textbf{The perception ``this is impermanent''} --- seeing change as it happens
\item \textbf{The perception ``this is not me''} --- seeing process instead of self
\end{enumerate}

Look at the structure. It's not random.\footnote{The Pali terms in order --- unwholesome: \textit{taṇhā, avijjā, lobha, dosa, moha, subhasaññā, sukhasaññā, niccasaññā, attasaññā}. Wholesome: \textit{samatha, vipassanā, alobha, adosa, amoha, asubhasaññā, dukkhasaññā, aniccasaññā, anattasaññā}. Note how the last four on each side are the \textit{vipallāsa} (distortions) and their corrections, corresponding to the four \textit{satipaṭṭhāna} (foundations of mindfulness). The text's summary verse makes this explicit: \textit{``Caturo ca vipallāsā, kilesabhūmī\ldots / Caturo satipaṭṭhānā, indriyabhūmī.''}} The first two on each side are the \textit{master pair}: craving and ignorance --- the two great drivers of suffering --- countered by calm and insight --- the two great liberators. Then come the three roots of unwholesome action --- greed, hatred, delusion --- mirrored by their three wholesome counterparts. And then come the four \textit{distortions of perception} --- seeing beauty, pleasure, permanence, and self where they don't ultimately exist --- mirrored by the four \textit{correct} perceptions.

The entire Canon, says the Netti, is a set of variations on the theme of replacing distorted perception with accurate perception. That's it. That's the whole game. Everything else --- the lists, the analyses, the stories, the rules --- is in service of this one cognitive revolution.

\begin{versebox}
Craving and ignorance, greed, hatred, delusion ---\\
four warped perceptions: the affliction's ground.
\stanza
Calm and insight, the three wholesome roots ---\\
four true perceptions: where liberation's found.
\stanza
Nine against nine, each paired with its twin ---\\
these eighteen roots: where all teaching begins.
\end{versebox}

\bigskip

That's your entire map. In the chapters that follow, we'll zoom in on each mode, each method, and see how they work on actual suttas. But right now, you already have the framework:

The Buddha's teaching is a system for replacing nine unwholesome roots with nine wholesome ones. The suttas are the delivery mechanism. The sixteen modes are how you read the delivery mechanism. The five methods are how you see the big picture. And the eighteen roots are what it all comes down to.

Simple. Elegant. And --- if the Netti is right --- complete.
